

\chapter{MATLAB Code Analysis: dlink.m}
\date{}

\section*{Overview}

\subsection*{What is the main purpose of this file?}
The main purpose of this script is to process a discrete half-edge data structure (referred to here as a ``double link'' structure) of a surface mesh, such as a Platonic solid. It derives topological relationships (edges, vertices, faces) and calculates geometric operators, ultimately aiming to build a discrete Laplace-Beltrami operator using ampli-twist (complex scaling and rotation) values. 

\subsection*{What type of analysis or calculation does it perform?}
It performs combinatorial and discrete geometric analysis. Specifically, it computes:
\begin{itemize}
    \item Topological orbits to identify vertices and faces from half-edge permutations.
    \item Angular defects at vertices (a discrete measure of Gaussian curvature).
    \item An incidence matrix connecting edges, vertices, and faces.
    \item A discrete Levi-Civita connection and a weighted cotangent Laplacian matrix.
\end{itemize}

\subsection*{What is the mathematical/theoretical context?}
The context lies in \textbf{Discrete Differential Geometry (DDG)} and \textbf{Topological Graph Theory}. It utilizes:
\begin{itemize}
    \item \textbf{Combinatorial Maps}: Using sets of ``darts'' or half-edges ($A$), permutation operators ($\phi$ for next half-edge around a face, $\theta$ for the opposite half-edge), and source mapping ($s$).
    \item \textbf{Discrete Connections}: The complex array $g = r \cdot e^{it}$ represents discrete holonomy or parallel transport (ampli-twist).
    \item \textbf{Homology/Cohomology}: Construction of the incidence boundary matrix $A$.
    \item \textbf{Discrete Laplace Operators}: The calculation of $L = W - \text{diag}(\sum W)$ mimics the cotangent Laplacian used on manifolds.
\end{itemize}

\section*{Step-by-Step Code Analysis}

\subsection*{Block 1: Variable Initialization and Involution Calculation}
\textbf{Explanation:} This section ensures an index array \texttt{idA} exists for the half-edges. It then computes the involution $\theta$, which pairs opposite half-edges. It finds the map such that $\theta(\phi(\theta(\phi))) = 1_A$ by matching the source nodes $s(\phi)$ and $s$. Next, it extracts the discrete topological features: vertices correspond to the orbits of $\theta$, and faces correspond to the orbits of $\phi$.
\begin{lstlisting}
try
idA=(1:nA)';
catch
idA=(1:numel(g))';
end

[~,theta_phi]=ismember([s(phi) s],[s,s(phi)],'rows');
theta(phi)=theta_phi; theta=theta(:);
[s_,sord]=orbits(theta); % the orbits of theta are the vertices
[q,qord]=orbits(phi); % the orbits of phi are the faces
idV=(1:max(s_))';  % identity on vertices
idF=(1:max(q))';
\end{lstlisting}

\subsection*{Block 2: Sorting and Ampli-Twist Extraction}
\textbf{Explanation:} This block creates sorted representations of vertices (\texttt{Vord}) and faces (\texttt{Ford}) alongside their original half-edge indices. It then extracts the geometric properties from the complex vector $g$: the amplitude (scaling factor $r$) and the twist (turning angle $t$ between adjacent half-edges). 
\begin{lstlisting}
Vord=sortrows([s_ sord, idA]);
Ford=sortrows([q qord, idA]);

g_=cumsum([0; exp(cumsum([0; log(g(Ford(1:end-1,3)))]))])

r=abs(g(Ford(:,3))); % ampli
t=angle(g(Vord(:,3))); % twist
\end{lstlisting}

\subsection*{Block 3: Angle Defect and Reference Angles}
\textbf{Explanation:} Here, the script calculates the angular defect at each vertex, which in discrete differential geometry is analogous to Gaussian curvature. It sums the turning angles $t$ around a vertex. It then computes reference angles (\texttt{refang}) to establish a local tangent space orientation for parallel transport.
\begin{lstlisting}
defang=full(cumsum(sparse(1,s(phi(Vord(:,3))),pi-t)))';
refang=2*pi*Vord(:,1).*cumsum(pi-t)./defang(Vord(:,1));
\end{lstlisting}

\subsection*{Block 4: Incidence Matrix Assembly}
\textbf{Explanation:} This section constructs a sparse block matrix $A$. This matrix acts as a discrete boundary or incidence operator. It is built such that if given functions defined on vertices ($fB$) and faces ($fC$), the linear system $A \backslash [fB; fC; \dots]$ interpolates or distributes these values onto the edges. 
\begin{lstlisting}
A=full([sparse(q,idA,1); sparse(s(phi),idA,1); sparse(orbits(theta(phi)),idA,1/2)]);
% matrix A is such that if given f0 on vertices and faces, say fB and fC with B the index set of vertices and C the index set of faces, we find f distributed by edges so that f0 is recoverd from f
% example 
fB=[1+1i; zeros(6,1); -1-1i]; fC=0*idF;
f=A\[fB; fC; zeros(numel(idA)/2,1)];
\end{lstlisting}

\subsection*{Block 5:   Levi-Civita Connection and Discrete Laplacian}
\textbf{Explanation:} The final part computes a discrete Levi-Civita connection using the reference angles to account for parallel transport across the curved discrete surface. It then defines a weight matrix $W$ (resembling the cotangent weight formula) and forces it to be symmetric. Finally, it constructs the graph Laplacian matrix $L$ as $L = W - D$, where $D$ is the degree (diagonal sum) matrix.
\begin{lstlisting}
% The connection:
Civitta=exp(1i*(refang - refang(theta(phi))) + pi);
% We could now redefine W as
W=sparse(s,s(phi),abs(g).*cot(angle(g)/2).*Civitta);
W=0.5*(W+W'); % force symmetry and get the equivalent 1/2(cot(a)+cot(b))
L=W-diag(sum(W));  %disp(full(L))
\end{lstlisting}

\section*{Expected Outputs and Visuals}

Based on the provided code, there are \textbf{no explicit visual plotting commands} (such as \texttt{plot}, \texttt{trisurf}, or \texttt{patch}) actively executed. 

\subsection*{Expected Console/Workspace Output:}
The primary output of this script is mathematical in nature. If executed, the MATLAB workspace will be populated with several structural arrays and sparse matrices representing the manifold:
\begin{itemize}
    \item \textbf{\texttt{A} (Incidence Matrix)}: A dense mapping matrix relating edges to faces and vertices.
    \item \textbf{\texttt{f} (Edge Distribution)}: A complex vector representing the distribution of a simulated function ($fB$, $fC$) across the edges.
    \item \textbf{\texttt{L} (Laplacian Matrix)}: A sparse, symmetric matrix representing the discrete Laplace-Beltrami operator. If the \texttt{disp(full(L))} command at the end of the script were uncommented, the console would print the dense numerical representation of this matrix.
\end{itemize}

\subsection*{Theoretical Visual Translation:}
If one were to visualize the results (e.g., using a custom rendering function later in the pipeline):
\begin{itemize}
    \item \textbf{Axes}: X, Y, and Z spatial coordinates for the vertices of the Platonic solid.
    \item \textbf{Plot Content}: A 3D wireframe or surface mesh.
    \item \textbf{Visualizing the Math}: The vector $f$ calculated from matrix $A$ could be represented as color gradients or directional arrows (vector fields) residing on the edges of the mesh. The angle defects (\texttt{defang}) could be plotted as color intensities at the vertices to show regions of high curvature.
\end{itemize}